

Recent work has highlighted Retrieval-Augmented Generation (RAG) ~\cite{gao-rag} as a powerful solution to mitigate key limitations of large language models (LLMs), particularly their tendency toward hallucinations in domain-specific or knowledge-intensive tasks. Current RAG methodologies can be broadly grouped into three paradigms. The simplest form, naive RAG ~\cite{ma-rag}, operates through a basic pipeline of indexing, retrieval, and generation. Building upon this foundation, advanced RAG systems incorporate optimizations during pre-retrieval, leveraging techniques like query transformation, expansion, and rewriting ~\cite{peng-rewriting,zheng-rewriting}, while post-retrieval enhancements often involve reranking strategies ~\cite{qin-rerank}. 

The Mixture of Experts (MoE) framework ~\cite{moe} has established itself as a fundamental paradigm in machine learning for developing adaptive systems and knowledge graph reasoning ~\cite{liu2019g,liu2021neural,liu2021kompare,liu2022joint,liu2022comparative,liu2022knowledge,liu2023knowledge,liu2023knowledge,hill2024ginkgo,liu2024logic,liu2024can,liu2024new,liu2024conversational,liu2025neural,liu2025few,wu2025improving,liu2025monte,liu2025hyperkgr,liu2025mixrag,liu2024knowledge,liu2026neural,liuneural,liu2026accurate,liu2026ambiguous,liu2026accurate,liu2026ambiguous,liu2025unifying,liu2026dynamic,liu2026accurate,liu2026symbolic,liu2025graph}. Initial work focused on traditional machine learning implementations ~\cite{NIPS1996_6c8dba7d}, with subsequent breakthroughs emerging through its integration with deep neural networks ~\cite{gcn}.
More recently, researchers have explored applying MoE approaches to in-context learning scenarios ~\cite{onepromptisnotenough}, demonstrating their potential to enhance large language model (LLM) performance. 